\subsolution{Дыбыс интерференциясы}{3.0}

Декарттық координаттар жүйесін енгізейік. Екі қабырғаның қиылысу сызығы $z$ осімен сәйкес келсін, ал қабырғалардың өздері $x=0$ және $y=0$ жазықтықтарында жатсын. Толқын $x \ge 0$, $y \ge 0$ квадрантында болады. Дыбыс толқындары бойлық болып табылады, сондықтан молекулалар толқындық вектордың бағытымен тербеледі. Түскен жазық толқын екі қабырғаға да $45^\circ$ бұрышпен координаттар басына қарай қозғалады. Оның $\vec{k}_1$ толқындық векторының компоненттері:
$$
k_x = k_y = k \cos(45^\circ) = \frac{k}{\sqrt{2}} = \frac{\sqrt{2}\pi}{\lambda}
$$
Әрі қарай \(k_0 = \pi \sqrt{2} / \lambda\) деп белгілейік. Онда $\vec{k}_1 = (-k_0, -k_0)$. Түскен толқынның жылдамдық векторы:

\begin{eq}[points=0.4]
    \vec{v}_1(x,y,t) = \left(-\frac{v}{\sqrt{2}}, -\frac{v}{\sqrt{2}}\right) \cos\left(\omega t + \frac{\sqrt{2}\pi}{\lambda} x + \frac{\sqrt{2}\pi}{\lambda} y\right),
\end{eq}

\criterion[type=remark]{\textbf{0.2 + 0.2 } вектордың дұрыс амплитудасы және фазаның координаттарға тәуелділігі үшін}

Дыбыс толқындары шағылысқанда фаза өзгермейді. Онда кеңістікте 3 түрлі шағылысқан толқынның суперпозициясы болады:

    \textbf{$x=0$ қабырғасынан шағылысқан толқын:} $\vec{k}_2 = (k_0, -k_0)$. Жылдамдықтың $x$-компонентасының таңбасы қарама-қарсыға өзгереді.

    \begin{eq}[points=0.1]
        \vec{v}_2(x,y,t) = \left(\frac{v}{\sqrt{2}}, -\frac{v}{\sqrt{2}}\right) \cos(\omega t - k_0 x + k_0 y),
    \end{eq}

    \textbf{$y=0$ қабырғасынан шағылысқан толқын:} $\vec{k}_3 = (-k_0, k_0)$. Жылдамдықтың $y$-компонентасының таңбасы қарама-қарсыға өзгереді.

    \begin{eq}[points=0.1]
        \vec{v}_3(x,y,t) = \left(-\frac{v}{\sqrt{2}}, \frac{v}{\sqrt{2}}\right) \cos(\omega t + k_0 x - k_0 y),
    \end{eq}

    \textbf{Екі қабырғадан да (бұрыштан) шағылысқан толқын:} $\vec{k}_4 = (k_0, k_0)$.

    \begin{eq}[points=0.1]
        \vec{v}_4(x,y,t) = \left(\frac{v}{\sqrt{2}}, \frac{v}{\sqrt{2}}\right) \cos(\omega t - k_0 x - k_0 y),
    \end{eq}

    \criterion[type=remark]{Жоғарыдағы балдар \textbf{тек} вектордың дұрыс амплитудасы \textbf{және} фазасы үшін беріледі}

Молекулалардың қорытқы жылдамдығын $\vec{v} = \vec{v}_1 + \vec{v}_2 + \vec{v}_3 + \vec{v}_4$ компоненттерді бөлек қосу арқылы табайық.
$x$-компонентасы үшін:
$$
v_x = \frac{v}{\sqrt{2}} \big[ -\cos(\omega t + k_0 x + k_0 y) + \cos(\omega t - k_0 x + k_0 y) - \cos(\omega t + k_0 x - k_0 y) + \cos(\omega t - k_0 x - k_0 y) \big]
$$

$y$ бойынша фазалары бірдей қосылғыштарды топтастырайық:
$$
v_x = \frac{v}{\sqrt{2}} \big[ \big(\cos(\omega t + k_0 y - k_0 x) - \cos(\omega t + k_0 y + k_0 x)\big) + \big(\cos(\omega t - k_0 y - k_0 x) - \cos(\omega t - k_0 y + k_0 x)\big) \big]
$$

Косинустар айырымының $\cos(A-B) - \cos(A+B) = 2\sin A \sin B$ формуласын қолданып:
$$
v_x = \frac{v}{\sqrt{2}} \big[ 2\sin(\omega t + k_0 y)\sin(k_0 x) + 2\sin(\omega t - k_0 y)\sin(k_0 x) \big]
$$

Ортақ көбейткішті жақша сыртына шығарып, синустар қосындысының $\sin(A+B) + \sin(A-B) = 2\sin A \cos B$ формуласын қолданамыз:
$$
v_x = \sqrt{2}v \sin(k_0 x) \big[ \sin(\omega t + k_0 y) + \sin(\omega t - k_0 y) \big] = 2\sqrt{2}v \sin(k_0 x) \cos(k_0 y) \sin(\omega t)
$$

Есептің симметриялылығына байланысты ($x \leftrightarrow y$ ауыстырғанда), жылдамдықтың $y$-компонентасы келесі түрде болады:
\begin{eq}[points=0.8, tail={, немесе жылдамдық компоненттерінің кем дегенде біреуі үшін баламалы қос тұрғын толқын теңдеуі үшін}]
    v_y = 2\sqrt{2}v \cos(k_0 x) \sin(k_0 y) \sin(\omega t),
\end{eq}

Көріп отырғанымыздай, жылдамдық компоненттері ``қос'' тұрғын толқын түрінде болады.

Молекулалар жылдамдығы модулінің квадраты $|\vec{v}|^2 = v_x^2 + v_y^2$:

\begin{align*}
    |\vec{v}|^2 & = 8v^2 \sin^2(\omega t) \big[ \sin^2(k_0 x) \cos^2(k_0 y) + \cos^2(k_0 x) \sin^2(k_0 y) \big] = 8v^2 \sin^2(\omega t) f(x,y)
\end{align*}

Жақша ішіндегі өрнектен екі түрлі координата бойынша туынды алсақ:

\begin{align*}
    \frac{\partial f(x,y)}{\partial x} & =  2k_0 \sin(k_0x) \cos(k_0 x) \cos^2(k_0 y) - 2 k_0 \cos(k_0 x ) \sin(k_0 x) \sin^2(k_0y)
    \\
    \frac{\partial f(x,y)}{\partial x} & = k_0 \sin(2 k_0 x) \cos(2 k_0y) = 0
\end{align*}

Дәл осындай шарт \(y\) бойынша туынды үшін де алынады. Экстремум үшін шарттарды аламыз:

\begin{gather*}
    \sin(2 k_0 x) \cos(2 k_0y) = 0, \quad  \sin(2 k_0 y) \cos(2 k_0x) = 0
    \\
     \left( x = \frac{\lambda}{2 \sqrt{2}} n \ \ \text{немесе} \ \  y =  \frac{\lambda}{2 \sqrt{2}} \left(\frac{1}{2} + m \right)\right) \ \ \text{және} \ \ \left( y = \frac{\lambda}{2 \sqrt{2}} n \ \ \text{немесе} \ \  x =  \frac{\lambda}{2 \sqrt{2}} \left(\frac{1}{2} + m \right)\right)
\end{gather*}

Экстремум шарттарынан \(f(x,y)\) функциясының максимумы \(1\)-ге тең екені түсінікті (мұны басқаша да көрсетуге болар еді). Онда бірден максимал жылдамдықты табуға болады:

\begin{eq}[points=0.5]
    v_{\max} = \sqrt{8v^2} = 2\sqrt{2}v,
\end{eq}

Максимум нүктелерінің координаттарын іріктеп алуға болады:
\begin{eq}[points=1.0]
    \left( \frac{\lambda}{2 \sqrt{2}} \cdot n, \frac{\lambda}{2 \sqrt{2}} \cdot m \right) \ \ \text{мұнда \(n,m\) сандарының тек біреуі тақ}.
\end{eq}

\criterion[type=remark]{\textbf{1.0} максимум координаттары үшін толық дұрыс және негізделген жауап үшін}

\criterion[type=remark]{\textbf{0.4} жартылай дұрыс жауап үшін, егер максимум нүктелері тікбұрышты \textbf{шексіз торда} болатыны көрініп тұрса}
